\documentclass{jsarticle}
\usepackage{amssymb}
\usepackage{amsthm}
\usepackage{amsmath}
\newtheorem{thm}{定理}
\newtheorem{prop}[thm]{命題}
\newtheorem{cor}[thm]{系}
\newtheorem{lemma}[thm]{補題}
\newtheorem{df}[thm]{定義}
\newtheorem{exam}[thm]{例}
\newtheorem{fact}[thm]{事実}
\renewcommand{\proofname}{証明}
\newcommand{\cc}{\mathbb{C}}
\title{複素解析と代数学の基本定理}
\author{電波通信}
\date{\today}
\begin{document}
\maketitle
この文書では複素解析を用いた代数学の基本定理の証明をする。複素解析を用いたこの基本定理の証明はリウビルの定理を用いる証明が簡潔で有名である。ここではリウビルの定理を用いる証明を含めて３つの証明を紹介する。証明に用いる3つの有名な事実をまず紹介する。

\begin{fact}[リウビルの定理]
$\cc$上で正則かつ有界な関数は定数関数に限る。
\end{fact}

\begin{fact}[最大絶対値の原理]
定数でない関数$f$は領域$D$上正則な関数な関数で$\bar{D}$上連続であるとする。この時$|f|$の最大値が存在するならば、内部では最大値を取り得ず、境界で最大値を達成する。
\end{fact}

\begin{fact}[弱い形のルーシェの定理]$D$を領域とし$\bar{\Delta}(a,R)\subset D$で$\gamma$を一回転しかしない$\Delta(a,R)$の境界（円周）とする。さらに$f,g$を$D$上の正則関数として$\gamma$上で
\[
|f(z)|>|f(z)-g(z)|
\]
が成り立ってるとする。このとき$\Delta(a,R)$に含まれる$f$と$g$の零点は同じ個数である。
\end{fact}
ルーシェの定理はもっと一般的な仮定のもとに述べられるが、この場では上の弱い形で十分である。

\begin{lemma}複素係数多項式の絶対値は$\cc$上最小値を持つ

\end{lemma}
\begin{proof}$\displaystyle \lim_{z\rightarrow \infty}|f(z)|=\infty$からわかる。
\end{proof}
この補題は多項式関数がリーマン球面に拡張出来ることの系である。

以下、$f$を定数でない複素係数多項式とする。

\begin{thm}[代数学の基本定理]定数でない複素係数多項式は$\cc$上零点を持つ
\end{thm}
\begin{proof}（リウビルの定理を用いる方法）$f$が零点を持たないとする。このとき
先の補題の証明から
\[
\lim_{z\rightarrow \infty}\frac{1}{f}=0
\]
がわかる。よって$\displaystyle \frac{1}{f}$は$\cc$上正則で有界であるから定数となるが、これは$f$が定数でないことに反するので$f$は零点を持つ。
\end{proof}

\begin{proof}（最大絶対値の原理を用いる方法）先の補題から最小値を達成する$a\in \cc$が少なくとも一つ存在する。さていま$f$は零点を持たないと仮定すると$f\neq 0$なので$\displaystyle \frac{1}{f}$も正則関数でこれに最大絶対値の原理を適用して$a$の十分小さい近傍の境界で$\displaystyle 
\frac{1}{f}$が最大値を取るものが存在するつまり
\[
\left|\frac{1}{f(a)}\right|<\left|\frac{1}{f(b)}\right|
\]
これは
\[
|f(b)|<|f(a)|
\]
を意味し$f(a)$が最小値ということに反する。よって$f$は零点を持つ。
\end{proof}

\begin{proof}（ルーシェの定理を用いる方法）$f(z)=z^n+a_{n-1}z^{n-1}\dots+a_1z+a_0$と書けると仮定しても良い。そして$g(z)=z^n$とする。このとき
\[
f(z)-g(z)=a_{n-1}z^{n-1}\dots+a_1z+a_0
\]
であり、これは高々$n-1$次多項式であり、$f$はn次多項式なので$R$を十分大きく取れば
$|z|=R$上で
\[
|f(z)>|f(z)-g(z)|
\]
が成り立ち、ルーシェの定理によって$f,g$は$|z|<R$上で同じ個数の零点をもつ。$g(z)=z^n$は$0$で$n$位の零点を持つので結局$f$は零点を持つ。
\end{proof}

\begin{thebibliography}{9}
\bibitem{1}Lars V.Ahlfors,Complex analysis,Third edition,McGraw-Hill,1979
\end{thebibliography}


\end{document}